% !TeX program = lualatex
% !TeX encoding = utf8
% !TeX root = ../Thermod.tex

%=========================================================
\Opensolutionfile{answer}[\currfilebase/\currfilebase-Answers]
\Writetofile{answer}{\protect\section*{\nameref*{\currfilebase}}}
\chapter{Теплота. Перший закон термодинаміки}\label{\currfilebase}
\makeatletter
\def\input@path{{\currfilebase/}}
\makeatother
%=========================================================






% --------------------------------------------------------
\section{1}
% --------------------------------------------------------

\paragraph{ }Якщо два шматка однакової маси, один залізний а інший свинцевий, обидва нагріті до 100\celsius, занурити у дві ізольовані посудини з водою однакової маси охолодженої до 0\celsius, то вичекавши поки в кожній із посудин не встановиться теплова рівновага, впевнимося, що посудина із шматком заліза покаже більшу температуру, ніж посудина зі свинцем. Навпаки, вода нагріта до 100\celsius, охолоджується шматком заліза більше, ніж шматком свинцю. Таким чином, необхідно робити відмінність між температурою та кількістю теплоти, при цьому за міру теплоти можна прийняти підвищення температури на один градус, що зазнає деяке тіло прийняте за еталон (вода, наприклад), приведене у контакт із досліджуваним тілом. При цьому припускають, що:
\begin{enumerate}
	\item [а)] кількість теплоти, яке відняли у охолоджуваного тіла дорівнює кількості теплоти одержаного охолоджувачем;
	\item [б)] інші причини зміни температури (стискання, тертя і таке інше) виключені.
\end{enumerate}

За одиницю виміру теплоти спочатку було прийнято, а також час від часу використовується і донині, <<калорія>> --- кількість тепла необхідна для нагріву 1 (одного) грама води від 14,5 до 15,5\celsius. Доречно порівняти походження слова «калорія» від латинського \emph{calor} (тепло) із \emph{temperatus} (спокійний, помірний – тобто такий, що може слугувати мірою рівноваги).

Здатність речовини сприяти або віддавати тепло можна схарактеризувати теплоємністю --- відношенням одержаного тепла до викликаної зміни температури:
\begin{equation}\label{2.1}
	C = \dfrac{\delta Q}{dt} = \dfrac{\delta Q}{dT} \left[\dfrac{\text{кал}}{\text{град}}\right]
\end{equation}

Зазначене визначення теплоємності кожен раз потребує уточнення того, яким чином відбувається передача тепла, тобто, як теплоємність, так і кількість теплоти залежить від способу передачі тепла або від процесу. У такому випадку (\ref{2.1}) належить писати:
\begin{equation}\label{2.2}
	C_x = \left(\dfrac{\delta Q}{dT}\right)_X
\end{equation}
де $X$ --- позначається процес (наприклад $P$, $V$, $H$, $M$...), а позначка $\delta Q$ вказує на те, що приріст теплоти не є повним диференціалом і завжди потребує знання процесу, при якому передавання тепла відбувається. Особливо легко це видно у випадку із газом. Насправді, газ можна нагріти на $\Delta T$ стискаючи його в адіабатичній оболонці, тобто взагалі не передаючи йому тепла. З другого боку на ту ж різницю калорій його можна нагріти підігріваючи полум'ям пальника та залишаючи його об’єм незмінним. Із наведеного прикладу очевидно, що значення теплоємності для однієї й тієї речовини може набувати всі значення від $-\infty$ до $+\infty$. Так, в адіабатичному процесі $C_s = \infty$ .

Джозеф Блек був перший, хто зумів експериментуючи з термометрами, встановити різницю між температурою та кількісною теплотою. Зокрема, він встановив, що в стані теплової рівноваги температури тіл, що знаходяться у термодинамічному контакті --- однакові, чим надзвичайно вразив сучасників. Ним же було введено поняття \emph{питомої теплоємності}, тобто теплоємності віднесеної до одиниці маси.
\begin{equation}\label{2.3}
	c = \dfrac{C}{m} \left[\dfrac{\text{кал}}{\text{град}\cdot\text{кг}}\right],\left[\dfrac{\text{кал}}{\text{град}\cdot\text{г}}\right]
\end{equation}

Тим самим, він спростував наявну на той момент думку про те, що теплоємність залежить тільки від маси речовини та не залежить від його складу. Ним же, разом із його учнем Джеймсом Уаттом було введено поняття про \emph{приховану теплоту} плавлення та випаровування. Замислившись одного разу над тим, чому кожну весну при таненні снігу ми не спостерігаємо нищівних паводків, він поставив експеримент, в якому порівняв час необхідний для підігріву води і льоду, що знаходяться при 0\celsius та розташовані на однаковій відстані від полум'я. Оскільки для нагріву льоду потребувалось часу у 4 рази більше ніж для нагрівання води на 1\celsius, Блек прийшов до висновку, що розтоплювання льоду потребує додаткового тепла назване їм прихованою теплотою або \emph{теплотою плавлення}.
\begin{equation}\label{2.4}
	\lambda = \dfrac{Q}{m}\left[\dfrac{\text{кал}}{\text{кг}}\right]
\end{equation}

% --------------------------------------------------------
\section{2}
% --------------------------------------------------------

Термодинаміка дає тільки співвідношення між фізичними величинами, але конкретні їх значення запозичують з досвіду або зі статистичної структури. В подальшому, скориставшись результатами молекулярно-кінетичної теорії, ми отримаємо вираз для теплоємності газів $C_\nu$ при сталому об'ємі:
\begin{equation}\label{2.5}
	C_\nu = \dfrac{i}{2} R \left[\dfrac{\text{кал}}{\text{моль}\cdot\text{град}}\right]
\end{equation}

Тут ідеться про \emph{молярну теплоємність}, чи теплоємність питому помножену на молярну масу:
\begin{equation}\label{2.6}
	C_\nu = c_\nu\cdot\mu\left[\dfrac{\text{кал}}{\text{моль}\cdot\text{град}}\right]
\end{equation}
$i$-кількість степенів свободи, що має молекула газу: $i=3$ для одноатомного газу, 5 для двохатомного та 6 для багатоатомного. Співвідношення (\ref{2.5}) справедливе в широкому діапазоні температур, але при більш високих температурах слід враховувати ефекти збудження та дисоціації молекул.

Систематичні досліди теплоємностей були проведені Дюлонгом і Пті. Вони з'ясували, що для більшості хімічно чистих елементів, що знаходяться в кристалічному стані, молярна теплоємність приблизно дорівнює $3R$ або близько \textbf{6} $\text{кал}/\text{град}\cdot\text{моль}$. Зараз відомо, що закон Дюлонга-Пті не виконується для речовин при температурах $T\ll\Theta_D$. При кімнатних температурах до них відносяться карбон, берилій та деякі інші. Для з'єднань Нейманом, Реньо, Каппом було встановлено, що в сполуці молярна теплоємність являє собою суму атомних теплоємностей, із якої складена сполука, причому незалежно від того, чи задовольняє закон Дюлонга-Пті чи ні властива атому теплоємність. Незважаючи на приблизний характер цієї закономірності для оцінки теплоємності сполук вона цілком придатна. Теплоємність рідин теоретично розраховується дуже важко. Якщо теплоємність твердого тіла біля точки затвердіння задовольняє закон Дюлонга-Пті, тоді теплоємність в рідкому стані має незначну відмінність від такої в твердому стані. Якщо ні, то в цьому випадку різниця особливо велика, так як це має місце у випадку з водою.

% --------------------------------------------------------
\section{3}
% --------------------------------------------------------

Дія будь якого \emph{калориметра} (приладу для визначення теплоємності) заснована на відніманні теплоти від одного тіла й передачі до другого та порівнянні результатів з еталонами, але вони нічого не дають нам в уявленні про природу теплоти. Більш того, схожість процесів теплообміну і встановлення теплової рівноваги з процесом перетікання рідини в сполучених посудинах, викликала уявлення про теплоту, як про деяку незнищенну субстанцію (\emph{теплород}), що перетікає від тіл з більш високої температурою до тіл з менш високою температурі при їх контакті --- теорії особливо популярної серед французької школи фізиків\footnote{Антон Лоран Лавуазьє, Жан Батіст Жозеф Фур'є, П'єр Сімон де Лаплас, Сімон Дені Пуассон, Саді Карно. Що стосується англійських природознавців, серед них із самого початку утвердилася думка, що природа теплоти якимось чином пов'язана із рухом (Ісаак Ньютон, Роберт Бойль)}.

Експериментальну перевірку теорії теплорода провів спочатку граф Румфорд (Бенджамін Томпсон). Займаючись свердленням гарматних стволів в зброярних військових майстернях в Мюнхені, він помітив, що тривалість виділення тепла залежить від тривалості свердління та виділяється його тим більше, чим тупіше свердло використовувалось та чим менше стружок утворювалось. Все це знаходилось в жахливій суперечності з наявними уявленнями про теплород, оскільки вважалось, що будь-який матеріал має кінцеву кількість теплорода, який здатен виділятися при руйнуванні матеріалу – отже чим більше руйнується матеріал, тим більше мало б виділятися тепла. Сам Румфорд дав таке пояснення своїм дослідам:<<\emph{Якщо ізольоване тіло чи система тіл здатні виробляти теплоту без обмежень, то остання не здатна бути матеріальною субстанцією; і мені вважається надзвичайно складним, якщо неможливим дати інше пояснення, крім того, що тільки рух здатен забезпечити збудження та поширення тепла в наших дослідах}>>\footnote{\emph{An Experimental Enquiry Concerning the Source of the Heat which is Excited by Friction}, Leslie, J. London. (1804)}. Згодом досліди Генрі Деві переконливо показали, що тертя здатне викликати плавлення льоду, яке, згідно з дослідами Д. Блека потребувало певної кількості тепла.

Всі ці досліди переконливо доводили, що механічна енергія завдяки тертю може бути вся перетворена в тепло. Навпаки, тепло може бути перетворено в роботу, наприклад, газ ізотермічно може виконувати роботу, піднімаючи вантаж. Ці міркування спочатку спонукали сера Румфорда, а потім Джеймса Прескота Джоуля оцінити кількість роботи необхідної для нагрівання речовини на одну калорію. Перший ще в 1978р.\footnote{<<Експериментальне дослідження природи теплоти, яку викликає тертя>> / An Experimental Enquiry Concerning the Source of the Heat which is Excited by Friction.} прикинув, що для виробництва 1 калорії потрібна робота, яка необхідна для підняття вантажу в 1кг на висоту $\sim5.5$м. В більш точних експериментах\footnote{Joule, J.P. (1843). "On the Calorific Effects of Magneto-Electricity, and on the Mechanical Value of Heat". Philosophical Magazine. 3. 23: 263, 347 \& 435. Джоуль розробив термометри, що вимірювали температуру з точністю до однієї двохсотої градуса.} Джоуль з точністю 0,5\% довів, що між кількістю теплоти й механічною енергією існує еквівалентність: 4,184 Дж роботи, незалежно від способу її перетворення, завжди виробляє 1 кал теплоти\footnote{Сучасне значення. Сам Джоуль знайшов в 1850 р. 4,159 Дж/кал}. Цікаво, що теоретична оцінка зроблена Юліусом Робертом фон Маєром в 1845р. 365 $kgf\cdot m/kcal$, а пізніше 425 $kgf\cdot m/kcal$ тривалий час не була відома широкому загалу фізиків\footnote{Die organische Bewegung im Zusammenhang mit dem Stoffwechsel (The Organic Movement in Connection with the Metabolism, 1845)}.

Вказану величину називають \emph{механічним еквівалентом теплоти}, а робота в СІ вимірюється в $\text{Дж}=\text{н}\cdot\text{м}$ саме на честь Джоуля. В гаусовій системі СГС роботу вимірюють в $\text{ерг} = 10^{-7}$ Дж, що в перекладі з грецької означає <<робота>>. В подальшому результати Джоуля перевірялись неодноразово. Таким чином, теплоту та роботу можна вважати проявленням однієї і тієї фізичної величини --- енергії. В перекладі з грецької $\varepsilon \nu \varepsilon \rho \gamma \varepsilon i\alpha$ означає одиниця ($\varepsilon\nu$ --- міра, запас) роботи, що начисто позбавляє це поняття містицизму. В цьому сенсі, якщо система володіє енергією, це означає, що вона здатна здійснювати роботу.

Енергія системи являє собою суму потенціальної та кінетичної енергії частинок, з яких вона складається. Якщо система являє собою поле, що взаємодіє з частинками, то слід враховувати енергію поля. По суті, за сучасним уявленням елементарні частинки (такі як електрони, кварки) це сильно збуджені стани полів таких, як електромагнітне, сильне відповідно. Поле в стані теплової рівноваги прийнято називати випромінюванням. Останнє також здатне змінювати енергію системи.

З точки зору термодинаміки енергія системи ділиться на зовнішню --- складову, що складається з енергії руху тіла та потенційної енергії системи в полі зовнішніх сил, та внутрішню --- всю іншу. \emph{Спосіб чи процес передачі енергії без зміни зовнішніх параметрів системи називається теплообміном, що виконується шляхом теплопровідності або радіацією, а енергію передану системі в такому процесі називається теплотою (або кількістю теплоти)}. Раніше вже зазначалось, що інший \emph{процес передачі енергії, зв'язаний зі зміною зовнішніх параметрів, називають відповідно, роботою}.

% --------------------------------------------------------
\section{4}
% --------------------------------------------------------

В попередніх розділах ми впевнились в тому, що енергія системи, зокрема внутрішня енергія, може бути змінена двояким способом: за допомогою передачі тепла, або здійсненням над нею роботи. Суттєво, що в термодинаміці постулюється, що внутрішня енергія являться функцією стану. Вперше, цей постулат, в формі близькій до сучасного звучання, був сформульований Юліусом Робертом фон Маєром\footnote{''On the Quantitative and Qualitative Determination of Forces'', 1841}. В теперішній час використовується декілька еквівалентних формулювань. Наведемо одне з них: <<\emph{При переході термодинамічної системи зі стану 1 в стан 2, отримана від навколишнього середовища сума робіт $\sum_{i}\delta A_i$ та теплоти $\delta Q$ визначається тільки станами 1 та 2. Ця сума не залежить від шляху переходу}>>. Звичайно ж \emph{ця сума дорівнює зміні внутрішньої енергії системи}. В математичному виразі:
\begin{equation}\label{2.7}
	dU = \delta Q - \sum_{i}\delta A_i
\end{equation}

Через те, що внутрішня енергія це функція стану, її приріст це повний диференціал та інтеграл від виразу (\ref{2.7}) по замкненому контуру перетворюється в нуль. Тоді для циклічних процесів $\oint dU = 0$ і як наслідок:
\begin{equation}\label{2.8}
	\delta Q = \sum_{i}\delta A_i
\end{equation}

Звідси можливо ще одне формулювання I закону термодинаміки, як \emph{неможливість існування perpetulum mobile I роду, тобто машини, що спроможна виконувати роботу без витрат тепла і яка за необхідністю повинна діяти циклічно} (за Максом Планком). По суті перший закон є не чим іншим, як законом збереження енергії поширеним на термодинамічні системи, оскільки очевидно, що для ізольованих систем $dU=0$, тобто $U=const$. Зрештою зауважимо, що твердження про те, що внутрішня енергія є функцією стану означає, що вона однозначно визначається внутрішніми (такі як T, P, ...) так і зовнішніми параметрами системи (V, H, ...).

% --------------------------------------------------------
\section{5}
% --------------------------------------------------------

Перший закон має безліч застосунків в термодинаміці, фізиці та хімії. Розглянемо перш за все співвідношення між найбільш важливими з точки зору практики теплоємностями $C_P$ та $C_V$. Згідно з (\ref{2.7}):
\begin{equation}\label{2.9}
	\delta Q = dU + \delta A,
\end{equation}
де для визначеності будемо вважати, що діє лише одне джерело роботи $\delta A = \mathbf{a} da$.
Тоді, вважаючи що $U$ це функція тільки температури $Т$ та зовнішнього параметра $a$, отримаємо:
\begin{equation}\label{2.10}
	C \equiv \dfrac{\delta Q}{\partial T} = \left(\left(\dfrac{\partial U}{\partial T}\right)_{a}dT + \left[\left(\dfrac{\partial U}{\partial a}\right)_{T} + \mathbf{a} \right]da\right)\dfrac{1}{dT}.
\end{equation}
Звідки
\begin{equation}\label{2.11}
	C_{a} = \left(\dfrac{\partial U}{\partial T}\right)_a,
\end{equation}
та
\begin{equation}\label{2.12}
	C_{\mathbf{a}} = \left(\dfrac{\partial U}{\partial T}\right)_{a} + \left[\left(\dfrac{\partial U}{\partial a}\right)_{T} + \mathbf{a}\right]\left(\dfrac{\partial a}{\partial T}\right)_{\mathbf{a}}.
\end{equation}
Зокрема, в окремому випадку $\mathbf{a}=P$, $a=V$, маємо
\begin{equation}\label{2.13}
	C_V = \left(\dfrac{\partial U}{\partial T}\right)_{V} ~~~~~~ C_P = \left[\left(\dfrac{\partial U}{\partial V}\right)_{T}+P\right]\left(\dfrac{\partial V}{\partial T}\right)_{P} + \left(\dfrac{\partial U}{\partial T}\right)_{V},
\end{equation}
та
\begin{equation}\label{2.14}
	C_P - C_V = \left[\left(\dfrac{\partial U}{\partial V}\right)_{T}+P\right]\left(\dfrac{\partial V}{\partial T}\right)_{P}.
\end{equation}
Зауважимо, що (\ref{2.14}) являє собою загальне співвідношення для будь-яких речовин, рівняння стану яких може бути записано як $P = P(V,T)$. Очевидно, що подібне співвідношення може бути отримано й для теплоємностей в інших процесах і для інших речовин, наприклад магнетика із рівнянням стану $M = M(H,T)$, звичайною заміною $P\leftrightarrow H$, а $V\leftrightarrow -M$ . З (\ref{2.14}) також стає очевидним, чому для твердих тіл та рідин не наголошується (чи наголошується але надзвичайно зрідка) відмінність в $C_P$ та $C_V$. Ця відмінність пропорційна коефіцієнту термічного розширення речовини, якій у твердих тілах та рідинах на декілька порядків менша, ніж у газів. Фізична причина різниці $C_P-C_V$ у газів зумовлена роботою проти зовнішнього тиску, що здійснюється газом при нагріванні.

Можна також зауважити, що виключаючи $U$ із (\ref{2.9}) можна перевірити теорію у наступний спосіб. Дійсно:
\begin{equation}\label{2.15}
	\left(\dfrac{\partial U}{\partial P}\right)_V =C_V\left(\dfrac{\partial T}{\partial P}\right)_V
\end{equation}

Із другого співвідношення (\ref{2.13}) знайдемо
\begin{equation}\label{2.16}
	\left(\dfrac{\partial U}{\partial V}\right)_{P} = (C_P - C_V)\left(\dfrac{\partial T}{\partial V}\right)_P -P
\end{equation}

Скориставшись тим, що $\dfrac{\partial^2 U}{\partial V\partial P} = \dfrac{\partial^2 U}{\partial P \partial V}$ знайдемо:
\begin{equation}\label{2.17}
	(C_P-C_V) \dfrac{\partial^2 T }{\partial P\partial V}+\left(\dfrac{\partial C_P }{\partial P}\right)_V\left(\dfrac{\partial T}{\partial V}\right)_P - \left(\dfrac{\partial C_V}{\partial V}\right)_P\left(\dfrac{\partial T}{\partial P}\right)_V=1
\end{equation}

Оскільки усі величини, що входять в (\ref{2.17}) доступні безпосередньому спостереженню, то з'являється можливість емпіричним шляхом перевірити справедливість першого закону.

% --------------------------------------------------------
\section{6}
% --------------------------------------------------------

З (\ref{2.13}) та (\ref{2.14}) очевидно, що для визначення $C_V$ досить знати лише калоричне рівняння стану $U = U(V,T)$, а для визначення $C_P$ слід знати ще й термічне рівняння стану. Правдиве й зворотне: із знання термічного рівняння стану, значень $C_V$ та $C_P$ інтегруванням можна встановити залежність $U = (V,T)$.

Як приклад зазначеного підходу розглянемо, як залежить від об'єму внутрішня енергія ідеального газу. Перші досліди були поставлені ще Гей-Люссаком та з точністю, що була досяжна в подібних експериментах давали незмінний результат: температура газу при його розширенні в вакуум не змінюється. Оскільки газ при розширенні роботи не виконував, теплотою з навколишнім середовищем не обмінювався, тоді згідно з першим законом $dU=0$. Оскільки $dU = \left(\dfrac{\partial U}{\partial T}\right)_VdT + \left(\dfrac{\partial U}{\partial V}\right)_TdV = C_VdT+\left(\dfrac{\partial U}{\partial V}\right)_T dV$, а зміна $dT=0$, то звідси слідує, що $\left(\dfrac{\partial U}{\partial V}\right)_T =0$, принаймні для газів близьких до ідеальних, з якими експериментував Гей-Люссак.

Більш точні виміри, проведені Томсоном разом з Джоулем (процес Джоуля-Томсона), здійснювали переведення газу необмеженим стаціонарним потоком із об'єму з великим тиском $P_1$ в простір з меншим тиском $P_2$, що досягається продавлюванням газу через трубку виготовлену з буксового дерева, та заткнутою в одному місці пористою перетинкою з очосу шовку чи вати. В стаціонарному стані для повітря виявилася мала, а для водню ледь вимірювана зміна температури газу і тому ми маємо право вважати, що для ідеального газу $\Delta T = 0$. Звідси витікає, що оскільки рівні кількості газу до та після перетинки займають різні об'єми їх внутрішня енергія залишається не змінною, а відтак не залежною від об'єму $U \neq U(V)$.

% --------------------------------------------------------
\section{7}
% --------------------------------------------------------

Для теоретичного осмислення отриманого результату уявимо поршень в положенні $ab$, котрий переміщуючись в положення $cd$ видавлює деяку кількість газу; останній проходячи перетинку видавлює інший поршень з положення $a'b'$ в положення $c'd'$ . Роботу, що здійснює перший поршень над газом, що міститься в об'ємі abcd буде: $P_1V_{abcd}=P_1V_1$, а робота, що була виконана газом: $P_2V_{a'b'c'd'}$. Оскільки за умовою теплообмін з навколишнім середовищем відсутній, тоді, згідно з першим законом:
\begin{equation}\label{2.18}
	U_2-U_2=(P_2V_2-P_1V_1)=0
\end{equation}
де $U_2$, $U_1$ --- внутрішня енергія газу після та до продавлювання. Звідси зауважуємо, що
\begin{equation}\label{2.19}
	U+PV=H=\text{const}
\end{equation}
в процесі Джоуля-Томсона. Величина H отримала назву \emph{ентальпія}, або тепловміст від грецької $\acute{\epsilon}\nu\theta\acute{\alpha}\lambda\pi\epsilon\i\nu$ (одиниця, запас, міра нагріву). Таким чином, процес Д-Т є ізоентальпійним.

Оскільки $U$ однозначно визначається $P$ та $V$, то ентальпія є функцією стану. Її приріст:
\begin{equation}\label{2.20}
	dH = dU + PdV +VdP = \delta Q+VdP
\end{equation}
в процесах, що проходять при постійному тиску дорівнює теплу переданому системі. Звідси отримаємо:
\begin{equation}\label{2.21}
	C_P = \left(\dfrac{\delta Q}{dT}\right)_P = \dfrac{dH}{dT}
\end{equation}

З отриманого виразу стає зрозумілим генезис слова тепловмісту. Оскільки з точки зору вимірювання теплоємності, легше виміряти $C_P$, а для твердих та рідких тіл $C_P\approx C_V$, то теплота отримана в процесі нагрівання була приблизно однакова при сталому об'ємі і тиску $\delta Q = C_P\Delta T\approx C_V\Delta T$ поводила себе як, висловлюючись сучасною мовою, функція стану, і тим самим переконуючи сучасників, що більша чи менша кількість теплоти знаходиться в тілі. Незважаючи на помилковість цих поглядів ентальпія отримала широке розповсюдження в техніці, так як дає безпосереднє уявлення про потік енергії і енергетичний баланс в відкритих системах. Так, уявимо замість поруватої перетинки деяке технічне обладнання (наприклад турбіну) через яке проходить робочий газ (водяна пара). Нехай ця машина виробляє корисну роботу (генерує електричний струм) $A$. І нехай, для загальності до неї було підведено (відведено) деяку кількість теплоти $Q$. Тоді рівняння балансу з урахуванням (\ref{2.19}) запишеться для одиниці маси газу:
\begin{equation}\label{2.22}
	h_1 +q = \dfrac{A}{m} +h_2
\end{equation}

Крім того в (\ref{2.22}) слід врахувати кінетичну енергію газу як цілого. Тоді:
\begin{equation}\label{2.23}
	h_1 +q = \dfrac{A}{m} +h_2 + \dfrac{v_2^2 - v_2^2}{2}
\end{equation}
де $v_2$ та $v_1$ швидкості стаціонарної течії газу в вихідному та вхідному перерізі пристрою. З (\ref{2.23}) легко отримати, що в окремих випадках при витікання газів із ракетного сопла, їх швидкість буде:
\begin{equation}\label{2.24}
	v_2 = \sqrt{2C_P\Delta T},
\end{equation}
де $\Delta T$ перепад температур всередині та ззовні сопла, а $v_1$ наближено дорівнює $0$.

% --------------------------------------------------------
\section{8}
% --------------------------------------------------------

Встановивши, що для ідеальних газів $\left(\dfrac{\partial U}{\partial V}\right)_T = 0$, можна отримати наступне важливе співвідношення:
\begin{equation}\label{2.25}
	C_P-C_V=P\left(\dfrac{\partial V}{\partial T}\right)_P =\dfrac{PV}{T} = R
\end{equation}
що отримало назву співвідношення Маєра. Сам Маєр отримав його використовуючи метод циклів, застосувавши його до процесу Гей-Люссака. Тоді для відрізків 13, 32, 21 відповідні зміни внутрішньої енергії складуть:
$$0; ~~ C_P(T_2-T_3) -P_2(V_3-V_2); ~~ C_V(T_1 - T_2)$$

Їх сума, згідно з першим законом дорівнює $0$ та, згідно з тим, що $T_1 = T_3$, $P_2V_3 = RT_3 = RT_1$, $P_2V_2 = RT_2$ отримаємо шукане співвідношення
\begin{equation}\label{2.26}
	C_P = C_V + R
\end{equation}
Зауважимо, що (\ref{2.25}) та (\ref{2.26}) дозволяють теоретично визначити механічний еквівалент тепла, що і було зроблено Робертом Маєром після спостережень над властивостями венозної та артеріальної крові.

% --------------------------------------------------------
\section{9}
% --------------------------------------------------------

Окрім ізотермічного, ізохоричного, ізобаричного процесу особливу в термодинаміці займають процеси політропічні, що відбуваються при постійному $C$. Для таких процесів, в загальному випадку, згідно із першим законом:
\begin{equation}\label{2.27}
	CdT = C_a dT + \left[\left(\dfrac{\partial U}{\partial a}\right)_T + \mathbf{a}\right]da
\end{equation}

Або використовуючи (\ref{2.12}):
\begin{equation}\label{2.28}
	(C-C_a) dT = \dfrac{C_{\mathbf{a}} - C_a}{(\partial a/\partial T)_{\mathbf{a}}}da
\end{equation}

З урахуванням того, що $dT = \left(\dfrac{\partial T}{\partial \mathbf{a}}\right)_a d\mathbf{a} + \left(\dfrac{\partial T}{\partial a}\right)_{\mathbf{a}} da$ отримаємо:
\begin{equation}\label{2.29}
	\left(\dfrac{\partial T}{\partial \mathbf{a}}\right)_a d\mathbf{a} + \dfrac{C_\mathbf{a} - C}{C_a -C}\left(\dfrac{\partial T}{\partial a}\right)_{\mathbf{a}} da =0
\end{equation}
Для окремого випадку адіабатичного процесу $C=0$. Тоді:\\
Приклад 1. Для $\mathbf{a} = p, a=V$
\begin{equation}\label{2.30}
	\left(\dfrac{\partial T}{\partial P}\right)_V dP+\dfrac{C_P}{C_V}\left(\dfrac{\partial T}{\partial V}\right)_PdV=0
\end{equation}

З урахуванням $PV = RT$ отримаємо
$$\dfrac{V}{R} dp =\gamma \dfrac{p}{R}dV =0, \text{ де } \gamma = \dfrac{C_P}{C_V}$$

Або інтегруючи маємо рівняння адіабати Пуассона:
\begin{equation}\label{2.31}
	PV^{\gamma} = \mathrm{const}
\end{equation}

З урахуванням рівняння стану
$$TV^{\gamma-1} = \mathrm{const}$$
Приклад 2. Для парамагнетиків $\mathbf{a}=H, a = M$ та $M = \dfrac{\chi}{T}H$.\\

Повторюючи викладки:
$$-\left(\dfrac{\partial T}{\partial H}\right)_M dH + \gamma \left(\dfrac{\partial T}{\partial M}\right)_H dM = 0$$
\begin{equation}\label{2.32}
	-\dfrac{\chi}{M} dH -\gamma\dfrac{\chi}{M^2}HdH=0=>\ln H=-\gamma \ln M
\end{equation}

Зрештою отримуємо рівняння адіабати для парамагнетиків $HM^{\gamma}= \mathrm{const}$

Існує безліч методів знаходження $\gamma$. Більшість з них використовують адіабатичність того чи іншого процесу. Прикладом може бути визначення $\gamma$ за швидкістю звуку, вирахуваної за формулою, що отримується в механіці.
\begin{equation}\label{2.33}
	C_S = \sqrt{\dfrac{\partial P}{\partial \rho}}
\end{equation}

Ньютон вважав процес поширення ізотермічним, тому отримав похідну під коренем $\dfrac{RT}{\mu}$. Насправді, внаслідок надзвичайно малої теплопровідності газів, стиснення та розширення в звуковій хвилі відбуваються не ізотермічно, а адіабатично. З урахуванням цього для вирахування (\ref{2.33}) слід використовувати (\ref{2.31}). Тоді:
$$\dfrac{\partial P}{\partial \rho} = \gamma \dfrac{RT}{\mu}$$

З урахуванням того що $C_S = 331.7~\dfrac{\text{м}}{\text{с}}$ при $T=273$\degree K отримаємо
$$\gamma =  \dfrac{\mu}{RT} \dfrac{\partial P}{\partial \rho} = \dfrac{28.9}{8.31\cdot10^7}\dfrac{(33170)^2}{273} = 1.40$$
що збігається із теорією.

% --------------------------------------------------------
\section{10}
% --------------------------------------------------------

Ще один принциповий спосіб вимірювання адіабатичної константи дає нам привід порівняти адіабати та ізотерми на індикаторній діаграмі. Передусім зауважимо, що через будь-яку точку на діаграмі ($P,V$) можна провести по одній адіабаті та одній ізотермі, що більш ніде не мають точок перетину (дотику). Для ідеальних газів в силу того, що $PV = \mathrm{const}$ для ізотерми та $PV^{\gamma} = \text{const}$ і адіабати це твердження є самоочевидним. Неможливість перетину двох ізотерм є наслідком нульового закону, згідно якого ми вважаємо температуру функцією стану, тобто однозначно визначеною через $P$ та $V$. Якщо ж точка перетину ізотерм існує, то виходячи із першої система може потрапити в стан з іншими $P$, $V$ та тією ж $T$. Що ж стосується до адіабат, то неможливість їх перетину є наслідком першого закону термодинаміки.

Звернемося тепер по другої частини твердження. Проведемо аргументи використовуючи методи циклів. Припустимо, що адіабата і ізотерма можуть перетинатися (наприклад т.1 і в т.2) утворимо замкнений цикл 1 І 2 ІІ між точками перетину. Тепло отримане на ізотермі 1 І 2 буде все перетворено в роботу, рівну площині заштрихованої фігури, що суперечить другому закону термодинаміки. Можна також показати, що будь-який інший політропічний процес що проходить через точку 0 перетину адіабати з ізотермою і лежить між ними, іде з від`ємною теплоємкістю. Для доказу слід скористатись (\ref{2.29}). В силу того, що для будь-якої речовини $\dfrac{C_P}{C_V}>1$ , нахил адіабати в точці перетину завжди крутіший, ніж нахил ізотерми: $\left(\dfrac{\partial P}{\partial V}\right)_S>\left(\dfrac{\partial P}{\partial V}\right)_T$. Для ідеального газу наприклад с урахуванням (\ref{2.31}) отримаємо:
\begin{equation}\label{2.34}
	\left(\dfrac{\partial P}{\partial V}\right)_T = -\dfrac{P}{V}
\end{equation}
\begin{equation}\label{2.35}
	\left(\dfrac{\partial P}{\partial V}\right)_S = -\gamma\dfrac{P}{V}
\end{equation}

З урахуванням того, що $\lpdr{P}{V}{S} = -\gamma \dfrac{\left(\partial T / \partial V\right)_P}{\left(\partial T / \partial P\right)_V}$ (що слідує із \eqref{2.30}) використовуючи правила трьох похідних знайдемо:
\begin{equation}\label{2.36}
	\lpdr{P}{V}{S} = \gamma\lpdr{P}{V}{T}
\end{equation}

Величину $K = -V\lppr{P}{V}$ називають \emph{модулем пружності}. З (\ref{2.36}) випливає, що адіабатичний та ізотермічний модулі пружності \underline{будь-якої} речовини, стан якого визначається змінними $T,P,V$, зв'язані між собою співвідношенням:
\begin{equation}\label{2.37}
	K_S = \gamma K_T
\end{equation}

Порівнюючи $K_S$ і $K_T$ можна знайти $\gamma$.

% --------------------------------------------------------
\section{11}
% --------------------------------------------------------

Ще одним наслідком першого закону, який за бажанням можна сформулювати і як сам перший закон, є \emph{закон Гесса}, широко застосовний в термохімії. Герман Гесс, проводячи серію дослідів по нейтралізації кислот, в 1840р сформулював наступне: ”Кількість тепла, що виділяється при утворенні будь-якої даної сполуки є сталим, не залежно від того чи утворилось ця сполука прямо чи непрямим шляхом, за один крок чи за серію кроків”. Дійсно, в реакції, що проходять при атмосферному тиску, енергія, що виділяється в процесі реакції, частково може бути перетворена в роботу. Тоді теплота, що виділяється (або поглинається) при постійному тиску є:
\begin{equation}\label{2.38}
	\begin{gathered}
		\Delta Q_P = \integ{U_1}{U_2}dU + \integ{V_1}{V_2}PdV =\\
	= (U_2 - U_1) + P(V_2-V_1) = (U_2 + P_2V_2) - (U_1 + P_1V_1) = H_2 -H_1 = \Delta H
	\end{gathered}
\end{equation}
де 1, 2 характеризує початковий та кінцевий стан і $P_1 = P_2 = P$. Теплота, яка виділяється дорівнює зміні ентальпії і її називають \emph{ентальпією} реакції. Реакція \underline{екзотермічна} якщо $\Delta H >0$, \underline{ендотермічна} при $\Delta H <0$.

Таким чином, загальне розв'язання задачі полягає в знанні ентальпії системи у всіх можливих її станах. Вартість такого підходу полягає в тому, що з'являється можливість визначення теплоти реакцій, які не можна виміряти прямо, або тих, яких важко виміряти.

Так, теплоти утворення окису вуглецю \ce{CO} із твердого вуглецю не може бути виміряна через те що вуглець ніколи не згорає до окису, а завжди із утворенням двоокису \ce{CO2}. Фавн і Зільберман знайшли тому .
\begin{equation}\label{2.39}
	\Delta H (\ce{C}_\text{Тв}+\ce{C}_\text{2газ} - \ce{CO}_\text{2газ}) = 97 \text{ ккал}/\text{моль}
\end{equation}

Теплоти згорання окису вуглецю з іншого боку
\begin{equation}\label{2.40}
	\Delta H (\ce{CO}_\text{газ}+\dfrac{1}{2}\ce{O}_\text{2газ} - \ce{CO}_\text{2}) = 68 \text{ ккал}/\text{моль}
\end{equation}

Віднявши друге з першого знаходимо:
\begin{equation}\label{2.41}
	\Delta H (\ce{C}_\text{Тв}+\dfrac{1}{2}\ce{O}_\text{2газ} - \ce{CO}_\text{газ}) = 29 \text{ ккал}/\text{моль}
\end{equation}

% --------------------------------------------------------
\section{12}
% --------------------------------------------------------

Теплота реакції, взагалі то кажучи залежить від температури, при котрій реакція проходить. Перший закон дає в цьому сенсі наступні вказівки. Нехай при деякій температурі $T$ теплоти реакції становлять $\Delta H$, а при $T+dT$ --- $\Delta H'$ відповідно.

Тоді:
\begin{equation}\label{2.42}
	Q' - Q = \Delta H' - \Delta H = (H_2'-H_1') - (H_2-H_1) = (H_2'-H_2) - (H_1'-H_1)
\end{equation}
Поділивши \eqref{2.42} на $dT$ і спрямувавши до нуля, отримаємо:
\begin{equation}
	\left(\dfrac{\delta Q}{\partial T}\right)_P = \lpdr{H}{T}{2} - \lpdr{H}{T}{1} = C_{P(2)} - C_{P(1)}
\end{equation}

Так, щоб визначити вплив температури при згоранні водню в рідину, слід порівняти теплоємність гримучого газу $\ce{H2} - \dfrac{1}{2} \ce{O2}$ із теплоємністю рідкої води. Молярна теплоємність гримучого газу дорівнює:
\begin{equation}\label{2.43}
	C_1 = 6.87 + \dfrac{1}{2} \cdot 7.04 = 10.29\text{ кал}/\text{моль}\cdot\text{град}
\end{equation}
а теплоємність води $C_2= 18 \text{ кал}/\text{моль}\cdot\text{град}$. Звідси
\begin{equation}\label{2.44}
	\dfrac{dH}{dT} = C_{P(2)} - C_{P(1)} = 7.61\text{ кал}/\text{моль}\cdot\text{град}
\end{equation}
І оскільки $\Delta H<0$ для цієї реакції, це означає, що теплота згоряння спадає з кожним градусом на $7.61$ кал.

% --------------------------------------------------------
\section{13}
% --------------------------------------------------------

Застосуємо перший закон до розгляду циклічних процесів. Історично склалось, що метод циклів був розроблений в технічній термодинаміці, але область його використання значно ширша оскільки будь-який періодичний процес допускає застосування цього методу. Почнемо розгляд з циклу Карно, котрий, внаслідок свого технічного значення зіграв велику роль в розвитку термодинаміки. Розберемо деякі властивості цього циклу, що працює на ідеальному газі, як робочій речовині. Використовуючи перший закон до кругового процесу знайдемо:
\begin{equation}\label{2.45}
	Q_{12} + Q_{23} + Q_{34} + Q_{41} + A_{12} + A_{23} + A_{34} + A_{41} =0
\end{equation}

На адіабатах $Q_{23} = Q_{41} =0$. Тоді робота на адіабатах дорівнює зміні внутрішньої енергії:
\begin{equation}\label{2.46}
	A_{23} = C_V(T_3-T_2) = -C_V(T_1-T_4)=-A_{41}
\end{equation}

Із (\ref{2.46}) слідує, що робота в циклі визначається тільки роботою на ізотермах:
\begin{equation}\label{2.47}
	A = A_{12} + A_{34} = \integ{V_1}{V_2}\dfrac{RT_1}{V}dV + \integ{V_3}{V_4} \dfrac{RT_3}{V}dV = R(T_1\ln\dfrac{V_2}{V_1} + T_3\ln\dfrac{V_4}{V_3})
\end{equation}
де на ізотермах $T_1 = T_2$ і $T_3 = T_4$. Стани 23 і 41 зв'язані рівнянням адіабати Пуассона адіабатичного процесу (\ref{2.31}). Звідси з урахуванням рівняння стану $PV^{\gamma} = \text{const} => TV^{\gamma-1} = \mathrm{const}$ отримаємо:
\begin{equation}\label{2.48}
	T_{1}V_{2}^{\gamma-1} = T_{3}V_{3}^{\gamma-1}; ~~~~ T_{3}V_{4}^{\gamma-1} = T_{1}V_{1}^{\gamma-1}
\end{equation}

Розділивши перше рівняння на друге отримаємо, що $V_2/V_1 = V_3/V_4$. Звідси із (\ref{2.47}) знайдемо, що
\begin{equation}\label{2.49}
	A = R(T_1-T_3) \ln\dfrac{V_2}{V_1}
\end{equation}
де, як і в (\ref{2.46}), \eqref{2.47} не обмежуючи спільність вважаємо, що робочої речовини є 1 моль.

На ізотермах в силу того, що $U_{12} = U_{34}=0$ маємо
\begin{equation}\label{2.50}
	Q_{12} =A_{12} = RT_1\ln\dfrac{V_2}{V_1} ~~~~ Q_{34} = A_{34} =- RT_3\ln\dfrac{V_2}{V_1}
\end{equation}

Із (\ref{2.49}), \eqref{2.50} робимо наступні корисні висновки. По-перше К.К.Д. такого циклу визначається, тільки відношенням температур:
\begin{equation}\label{2.51}
	\eta = \dfrac{A}{Q_{12}} = \dfrac{R(T_3-T_1)\ln \dfrac{V_2}{V_1}}{RT_1\ln \dfrac{V_2}{V_1}}=1-\dfrac{T_3}{T_1}
\end{equation}
По-друге: між величинами $Q_{12}$, $Q_{34}$ та $A= A_{12} + A_{34}$ існує співвідношення:
\begin{equation}\label{2.52}
	Q_{12} : Q_{34} : A = -T_1:T_3:(T_1-T_3)
\end{equation}
яке еквівалентне наступному:
\begin{equation}\label{2.53}
	\dfrac{Q_{12}}{T_1} +\dfrac{Q_{34}}{T_3} =0
\end{equation}

% --------------------------------------------------------
\section{14}
% --------------------------------------------------------